\section{Conclusion et perspectives}

\subsection{Bilan}

La question posée en introduction était simple à énoncer : la turbulence ambiante
homogène et isotrope accélère-t-elle ou ralentit-elle la vitesse d'ascension d'une
bulle de gaz isolée, à nombre de Bond $\mathrm{Bo}=1$ et nombre de Galilée
$\mathrm{Ga}=70$ fixés (donc au-delà du seuil d'instabilité de trajectoire
$\mathrm{Ga}_c\simeq44$, régime de trajectoire oblique) ? La réponse, obtenue par
simulation numérique directe avec Basilisk, est claire : \textbf{la
turbulence ralentit la remontée}, et le ralentissement relatif croît avec
l'intensité turbulente relative $\beta=u'/u_\infty$, le paramètre de contrôle mis
en avant par la revue de Legendre \& Zenit \cite{legendre2025}.

\begin{table}[H]
    \centering
    \renewcommand{\arraystretch}{1.3}
    \small
    \begin{tabular}{lcccccccccc}
        \toprule
        $\beta$ & $0{,}05$ & $0{,}075$ & $0{,}10$ & $0{,}15$ & $0{,}22$ & $0{,}31$ & $0{,}33$ & $0{,}38$ & $0{,}50$ & $0{,}65$ \\
        \midrule
        $u_{\infty,\mathrm{turb}}/u_\infty - 1$ & $-3\%$ & $0\%$ & $-4\%$ & $-6\%$ & $-12\%$ & $-31\%$ & $-32\%$ & $-38\%$ & $-58\%$ & $-72\%$ \\
        \bottomrule
    \end{tabular}
    \caption{Réduction relative de la vitesse d'ascension mesurée en fonction de
    l'intensité turbulente relative $\beta$ ($n=5$ réalisations par point, 10 points acquis).}
    \label{tab:bilan_beta}
\end{table}

Ces dix points (tab.~\ref{tab:bilan_beta}) suivent une loi quadratique
$1-c\beta^2$ aux plus faibles intensités ($\beta \le 0{,}15$) et s'infléchissent aux plus fortes, un
comportement qui \textbf{rejoint, sans aucun ajustement de paramètre, la courbe
de référence de Liu, Farsoiya, Perrard \& Deike \cite{liu2024}} : transition entre le
régime parabolique à petit nombre de Froude turbulent $\mathrm{Fr}'=u'/\sqrt{gD}$
et le régime asymptotique en $0{,}37/\mathrm{Fr}'$ aux forts Froude ($\beta \ge 0{,}50$). Nos points capturent
cette transition et confirment la branche asymptotique. Le mécanisme
physique responsable de ce ralentissement a été identifié comme étant la
\textbf{traînée non linéaire} plutôt que l'échantillonnage préférentiel :
l'échantillonnage préférentiel, mesuré directement, reste faible et
essentiellement nul en moyenne ($\le12\%$), alors que les fluctuations de la
vitesse de glissement bulle-fluide, elles, croissent nettement avec $\beta$ (de
$8\%$ à $50\%$) et suffisent, via la non-linéarité de la loi de traînée, à
expliquer le ralentissement observé. Ce résultat est cohérent avec le régime
identifié par Liu \emph{et al.} \cite{liu2024} pour une bulle de taille
supérieure à l'échelle de Kolmogorov, qui est le régime du présent
travail.

Ce résultat central n'a de valeur que parce que la chaîne de calcul qui le
produit a été validée de façon quantitative. La référence en liquide au repos
donne une vitesse terminale $u_\infty=12{,}27$, convergée en maillage à
$0{,}32\%$ près entre les niveaux 7 et 8 ($10{,}93$ / $12{,}27$ / $12{,}31$ aux
niveaux 6/7/8) ; les courants parasites représentent $0{,}08\%$ de $u_\infty$ ;
le biais lié au sillage périodique est borné et vaut $-0{,}8\%$ ; le volume de la
bulle reste conservé entre $99{,}05\%$ et $99{,}88\%$ tout au long des
simulations ; et la résolution turbulente satisfait partout $k_{\max}\eta\ge3{,}92$,
très au-dessus du seuil de résolution $1{,}5$. Le résultat central, à savoir le
ralentissement de $-3\%$ à $-72\%$ et son accord parfait avec la courbe de référence de
\cite{liu2024}, repose donc sur une chaîne numérique dont chaque maillon a été
quantifié et non simplement supposé correct.

\subsection{Leçon de méthode}

Le résultat le plus utile de ce stage n'est peut-être pas le nombre final, mais
la façon dont il a failli être faux. Une version antérieure du code, utilisant le
module standard de gravité réduite (\texttt{reduced.h}), produisait un résultat
en apparence tout aussi convaincant : une bulle qui montait, ralentissait sous
l'effet de la turbulence, puis semblait ne plus atteindre de vitesse terminale
nette. Ce résultat était spectaculaire, cohérent avec une partie de
l'intuition physique, et \textbf{entièrement faux} : la gravité réduite exprime
la poussée par un potentiel linéaire en $z$, donc \emph{non périodique}, ce qui
est incompatible avec un domaine périodique selon la direction de la gravité,
la bulle se retrouvait purement et simplement épinglée à la frontière du
domaine par une force parasite, quelle que soit la turbulence appliquée.

Ce qui a permis de démasquer cette erreur numérique n'est pas une relecture attentive du
code, mais un \textbf{cas limite dont la réponse est connue a priori} : une
bulle lâchée dans un liquide au repos \emph{doit} monter
indéfiniment à sa vitesse terminale, sans jamais s'arrêter. Faire tourner ce cas
laminaire, en apparence le moins intéressant de la campagne, a montré que la
bulle s'y comportait exactement comme dans les cas turbulents, elle
s'immobilisait à la même hauteur, ce qui est physiquement impossible en
l'absence de toute turbulence et a immédiatement désigné le mécanisme en cause.
Sans ce test, j'aurais présenté comme résultat de stage l'énoncé « la turbulence
supprime la remontée de la bulle », une conclusion plausible, alignée avec une
partie de la littérature, et qui aurait été une pure erreur numérique. C'est le
fil rouge méthodologique de ce travail : dans une simulation multi-physique où
plusieurs effets se combinent, un cas limite au comportement connu à l'avance
est le test de non-régression le plus informatif qu'on puisse construire,
souvent bien plus que le cas « intéressant » que l'on cherche à produire.

\subsection{Limites}

Ce travail comporte plusieurs limites qu'il convient d'assumer explicitement.

\begin{itemize}
    \item Un seul couple de paramètres $(\mathrm{Bo},\mathrm{Ga})=(1,70)$ a été
    exploré : rien ne garantit que le comportement observé, ni les valeurs de
    la loi de réduction, se transpose sans changement à d'autres couples,
    notamment à des bulles plus ou moins déformables.
    \item Les dix points acquis couvrent $\beta$ de $0{,}05$ jusqu'à $0{,}65$ (les deux derniers points de la campagne, $\beta=0{,}85$ et $\beta=1{,}00$, sont en attente de calcul sur le cluster).
    \item Les barres d'erreur croissent avec $\beta$ (jusqu'à $\pm6{,}7\%$ à $\beta=0{,}65$), chaque point reposant sur une moyenne d'ensemble de $n=5$ réalisations turbulentes indépendantes.
    \item Le domaine est triplement périodique : la bulle finit par recroiser
    son propre sillage à travers les conditions périodiques, ce qui introduit un
    biais chiffré plus haut ($-0{,}8\%$) mais non nul, une bulle en domaine
    réellement infini ne serait pas exposée à cet effet.
    \item La mesure du mécanisme (traînée non linéaire contre échantillonnage
    préférentiel) reste approximative : elle repose sur une coquille définie à
    l'échelle intégrale de la turbulence autour de la bulle, un choix
    raisonnable mais non unique, et non sur une décomposition rigoureuse des
    contributions au bilan de quantité de mouvement.
\end{itemize}

\subsection{Perspectives}
\label{sec:perspectives}

Plusieurs prolongements directs de ce travail sont identifiés.

\begin{itemize}
    \item \textbf{Étendre le balayage en $\beta$ jusqu'à $\beta\simeq1$.} Une
    campagne est en cours, qui triple la plage explorée par rapport aux cinq
    points actuels ; elle est rendue possible par le fait qu'à $g=4$ le nombre de
    Weber turbulent $\mathrm{We}_t$ reste sous le seuil de fragmentation
    $\mathrm{We}_c\simeq3$ mesuré par Rivière \emph{et al.} \cite{riviere2021},
    ce qui garantit que la bulle reste intacte sur toute la plage visée.
    \item \textbf{Tester plus finement la loi quadratique à petit $\beta$}, en
    resserrant l'échantillonnage aux faibles intensités turbulentes pour
    contraindre la constante $c$ de la loi $1-c\beta^2$ avec moins
    d'incertitude.
    \item \textbf{Explorer la déformabilité de la bulle}, en faisant varier le
    nombre de Bond $\mathrm{Bo}$ à $\mathrm{Ga}$ fixé, pour tester la
    généralité du résultat au-delà du couple unique traité ici.
    \item \textbf{Découpler $\beta$ et $\mathrm{Re}_\lambda$}, actuellement liés
    par construction dans nos précurseurs, pour isoler l'effet propre de
    l'intensité turbulente relative de celui du nombre de Reynolds turbulent.
    \item \textbf{Affiner la mesure du mécanisme} en multipliant les
    réalisations indépendantes par point de $\beta$, ce qui réduirait le bruit
    statistique sur la décomposition échantillonnage préférentiel / traînée non
    linéaire et sur les barres d'erreur identifiées ci-dessus comme la
    principale limite du résultat actuel.
\end{itemize}
